\section{Introduction}
\label{sec:intro}

Human review is often justified by asking whether the human knows something the model does not.
In clinical AI, a radiologist may catch findings that an imaging model misses, and residual-expertise audits can test whether the human signal carries predictive information beyond the model \citep{AlurEtAl2023,AlurRaghavanShah2024}.
But deployment rewards actions, not predictive information.
A radiologist can shift the posterior without changing the admit/discharge decision; only posterior movements that cross the reward-induced action boundary matter in practice.
Residual predictive expertise and residual decision value are not the same estimand.

Complementarity is therefore reward-relative.
The same human signal can matter under one reward matrix and not under another, not because the human's expertise changed, but because the boundary moved.
Human--AI complementarity is not a property of the prediction problem alone.
For deployment, it also depends on the model belief, the human signal, the action set, and the reward matrix.
When audits treat residual predictive expertise as residual decision value, they can spend scarce review budget on cases where the human is informative but operationally irrelevant.

We formalize the missing estimand as \emph{boundary regret}: the regret of retaining the model-only action after observing the human signal.
For model belief $b_x$, model-only action $a_x = \argmax_a r_a \cdot b_x$ (optimal under $b_x$), human-updated belief $b_{x,h} = P(Y \mid x, h)$, and terminal value $V(b) = \max_a r_a \cdot b$,
\begin{equation}
  \BR(x, h) \;=\; V(b_{x,h}) \;-\; r_{a_x} \cdot b_{x,h},
\end{equation}
which is positive exactly when a rival action becomes preferred under $b_{x,h}$.
In any finite-action Bayesian decision problem, the value of observing the human signal equals expected boundary regret (Theorem~\ref{thm:voi-br}):
\begin{equation}
  \VoI(H \mid b_x) \;=\; \E_{h \mid x}\bigl[\BR(x, h)\bigr].
\end{equation}
Predictive gain without a reachable boundary crossing has zero decision value for the deployed action set.
We formalize in Lean~4 with Mathlib \citep{DeMouraUllrich2021,MathlibCommunity2020} both the zero-regret boundary condition and the value-of-information identity above.
The paper develops this criterion in four steps:
\begin{itemize}[nosep,leftmargin=1.5em]
  \item We define boundary regret and prove, in Lean~4, that $\VoI(H \mid b_x) = \E_{h \mid x}[\BR(x,h)]$, with a geometric facet-activation criterion that characterizes when boundary regret is strictly positive (Theorems~\ref{thm:voi-br}--\ref{thm:facet-bridge}).
  \item We validate the plug-in estimator $\BRhat$ on synthetic finite-action problems, including a strong-null condition where humans carry substantial predictive signal yet boundary regret is zero everywhere: predictive gain is not sufficient for decision value.
  \item We show on CheXpert that residual predictive expertise and residual decision value separate as a function of reward geometry, exactly as the theory predicts: the same model, the same readers, and the same conditions yield opposite-signed correlations under different rewards.
  \item We show that ex ante $\BRhat$ improves scarce-review allocation over reward-aware margin, reward-agnostic uncertainty, residual-expertise, and learning-to-defer baselines where boundary-regret mass is reachable, and converges to reward-aware margin where it is not (Theorem~\ref{thm:facet-bridge}).
\end{itemize}
